
En el contexto de datacenters, las VPNs (\textit{Virtual Private Networks}) se utilizan principalmente como mecanismo de interconexión segura entre sedes, no como solución de acceso remoto para usuarios individuales. Su rol es garantizar que el tráfico que atraviesa redes compartidas o no confiables como Internet o redes de operador, llegue cifrado y autenticado al destino, complementando el aislamiento lógico que ofrecen los protocolos de overlay.

\subsection{IPsec como base de las VPNs site-to-site}

IPsec (\textit{Internet Protocol Security}) es el conjunto de protocolos más ampliamente utilizado para implementar VPNs site-to-site en entornos de datacenter. Opera en capa de red y proporciona tres servicios fundamentales: autenticación de origen, integridad de los datos y cifrado del contenido. Su modo de operación más utilizado para interconexión de sedes es el \textit{modo túnel}, en el que el paquete IP original se encapsula completamente dentro de un nuevo paquete IPsec, ocultando las direcciones IP internas durante el tránsito por la red pública.

IPsec se apoya en dos protocolos complementarios: AH (\textit{Authentication Header}), que provee autenticación e integridad sin cifrado, y ESP (\textit{Encapsulating Security Payload}), que añade cifrado al tráfico encapsulado y es el componente utilizado en la mayoría de las implementaciones de VPN. La negociación de parámetros de seguridad —algoritmos, claves y políticas— se gestiona mediante IKE (\textit{Internet Key Exchange}), que establece las asociaciones de seguridad (SA) entre los extremos del túnel.

\subsection{GRE sobre IPsec}

Una combinación frecuente en redes empresariales y de datacenter es GRE sobre IPsec: GRE proporciona el mecanismo de tunelización y soporte para múltiples protocolos, mientras que IPsec agrega la capa de cifrado y autenticación que GRE no ofrece por sí solo. En esta arquitectura, el paquete original se encapsula primero con GRE y luego el resultado se cifra con IPsec en modo túnel, permitiendo transportar tráfico de forma segura entre sedes a través de redes no confiables.

\subsection{MPLS VPN: solución a nivel de operador}

En entornos donde la interconexión de datacenters se realiza a través de la red de un operador de telecomunicaciones, MPLS ofrece una alternativa a las VPNs basadas en IPsec. Dado que los paquetes van a viajar dentro de la red del operador, no es necesario cifrar el contenido, solo basta con aíslar lógicamente. Esto resulta una ventaja en lo que respecta a escalabilidad y calidad de servicio.

\subsection{Rol de las VPNs en la arquitectura de overlay}

Las VPNs no reemplazan a los protocolos de overlay sino que los complementan: mientras VXLAN, GRE o EVPN-VXLAN proveen el plano de datos y el aislamiento lógico entre redes virtuales, las VPNs site-to-site basadas en IPsec o MPLS garantizan que ese tráfico viaje de forma segura cuando debe atravesar infraestructura compartida o no controlada. Esta distinción es fundamental: el aislamiento lógico evita interferencias entre tenants, pero solo el cifrado garantiza confidencialidad frente a observación o interceptación del tráfico en tránsito.
